\documentclass[12pt,a4paper]{article}
\usepackage[utf8]{inputenc}
\usepackage[T1]{fontenc}
\usepackage{geometry}
\geometry{margin=1in, headheight=15pt}
\usepackage{listings}
\usepackage{xcolor}
\usepackage{hyperref}
\usepackage{titlesec}
\usepackage{enumitem}
\usepackage{booktabs}
\usepackage{graphicx}
\usepackage{fancyhdr}
\usepackage{tcolorbox}

\definecolor{codegreen}{rgb}{0,0.6,0}
\definecolor{codegray}{rgb}{0.5,0.5,0.5}
\definecolor{codepurple}{rgb}{0.58,0,0.82}
\definecolor{backcolour}{rgb}{0.95,0.95,0.92}

\lstdefinestyle{mystyle}{
    backgroundcolor=\color{backcolour},
    commentstyle=\color{codegreen},
    keywordstyle=\color{magenta},
    numberstyle=\tiny\color{codegray},
    stringstyle=\color{codepurple},
    basicstyle=\ttfamily\footnotesize,
    breakatwhitespace=false,
    breaklines=true,
    captionpos=b,
    keepspaces=true,
    numbers=left,
    numbersep=5pt,
    showspaces=false,
    showstringspaces=false,
    showtabs=false,
    tabsize=2
}
\lstset{style=mystyle}

\pagestyle{fancy}
\fancyhf{}
\rhead{Monopoly-LK}
\lhead{Function Documentation}
\cfoot{\thepage}

\hypersetup{
    colorlinks=true,
    linkcolor=blue,
    filecolor=magenta,
    urlcolor=cyan,
}

\title{
    \textbf{Monopoly-LK Simulation} \\
    \large Complete Function Logic and Documentation \\
    \vspace{0.5cm}
    \small SCS1301 - Programming Fundamentals
}
\author{}
\date{\today}

\begin{document}

\maketitle
\tableofcontents
\newpage

% ============================================================
% SECTION 1: DATA STRUCTURES
% ============================================================
\section{Data Structures and Type Definitions (\texttt{types.h})}

The program defines several \texttt{struct} and \texttt{enum} types that form the foundation of the simulation.

\begin{tcolorbox}[colback=gray!10, colframe=gray!60, title=\textbf{Key Data Structures}]
\begin{itemize}[noitemsep]
    \item \textbf{player} --- Holds all state for a player: cash, position, properties, loans, jail status, etc.
    \item \textbf{square} --- Represents one board square (property, railway, utility, event, tax, etc.) with nested structs for each type's specific attributes.
    \item \textbf{context} --- Tracks global game state: current round, active economic event, inflation rate, and interest rate.
    \item \textbf{playerPointers} --- Aggregates pointers to all four player objects plus the bank for convenient passing.
    \item \textbf{diceRollType} --- Stores a dice roll result and whether doubles were rolled.
\end{itemize}
\end{tcolorbox}

Player types are defined as an enum:
\begin{lstlisting}[language=C]
typedef enum {
    noOwnner = 301,
    bankOfCeylon,
    aggresiveInvester,
    conservativeBanker,
    riskTaker,
    opportunisticTrader
} playerType;
\end{lstlisting}

Square types categorize each board position:
\begin{lstlisting}[language=C]
typedef enum {
    go, special, railway, utility, event,
    insure, tax, bank, property
} squareType;
\end{lstlisting}

\newpage

% ============================================================
% SECTION 2: ENTRY POINT
% ============================================================
\section{Entry Point (\texttt{main.c})}

\subsection{\texttt{main(void)}}
The entry point is minimal; it simply calls \texttt{startGame()}.
\begin{lstlisting}[language=C]
int main(void) {
    startGame();
    return 0;
}
\end{lstlisting}

\textbf{Logic:}
\begin{enumerate}[noitemsep]
    \item Call \texttt{startGame()} to initialize and run the entire simulation.
    \item Return 0 on success.
\end{enumerate}

\newpage

% ============================================================
% SECTION 3: GAME INITIALIZATION AND MAIN LOOP
% ============================================================
\section{Game Initialization and Main Loop (\texttt{game.c})}

\subsection{\texttt{startGame(void)}}
This is the central function that initializes the game, runs the main loop, and determines the winner.

\textbf{Step-by-step logic:}
\begin{enumerate}[noitemsep]
    \item \textbf{Initialize context:} Set all context fields (inflation, economic event, round counter, interest rate) to zero or -1.
    \item \textbf{Initialize card deck indices:} Set \texttt{topNationalEventCard} and \texttt{topreigionaldevelopmentcard} to 0.
    \item \textbf{Create players:} Instantiate four players (Aggressive Investor, Conservative Banker, Risk Taker, Opportunistic Trader), each starting at square 1 with LKR 30,000.
    \item \textbf{Create bank:} Instantiate the \texttt{Bank of Ceylon} as a player entity.
    \item \textbf{Set up player pointers:} Create a \texttt{playerPointers} struct linking all player objects.
    \item \textbf{Randomize seed:} Call \texttt{srand(time(NULL))} for randomness.
    \item \textbf{Determine turn order:} Call \texttt{ranker()} to roll dice and rank players, then \texttt{finalRankAssign()} to assign ranked players to positions 1--4.
    \item \textbf{Initialize board:} Call \texttt{initializeBoard(board)} to populate all 40 squares.
    \item \textbf{Print game start info:} Display player names, starting cash, and turn order.
    \item \textbf{Main game loop:} While \texttt{true}:
    \begin{enumerate}[noitemsep]
        \item Break if round exceeds 500.
        \item For each non-bankrupt player, call \texttt{move()}.
        \item Call \texttt{networthEvaluate()} to recalculate all net worths.
        \item Call \texttt{bankruptCheck()} to mark bankrupt players; break if 3 are bankrupt.
        \item Call \texttt{roundCounter()} to advance the round and print summaries.
        \item Every 15 rounds, trigger \texttt{econEventActivate()}.
        \item Every 10 rounds, trigger \texttt{inflationRateRelease()}.
    \end{enumerate}
    \item \textbf{Determine winner:} If 3 players are bankrupt, the remaining player wins. Otherwise, the player with the highest net worth wins.
    \item \textbf{Announce winner:} Print the winner's name and net worth.
\end{enumerate}

\newpage

% ============================================================
% SECTION 4: BOARD INITIALIZATION
% ============================================================
\section{Board Initialization (\texttt{board.c})}

\subsection{\texttt{initializeBoard(square *board)}}
Populates the \texttt{board[40]} array with all square data.

\textbf{Step-by-step logic:}
\begin{enumerate}[noitemsep]
    \item Create a \texttt{bankOfCeylonInstitution} as the default owner.
    \item For each of the 40 squares (0--39), assign:
    \begin{itemize}[noitemsep]
        \item \textbf{squareID}: Index 0--39.
        \item \textbf{name}: Location name (e.g., ``Pettah'', ``Colombo Fort Railway Station'').
        \item \textbf{type}: One of \texttt{go, property, railway, utility, event, insure, tax, bank, special}.
        \item \textbf{owner}: Defaults to \texttt{\&bankOfCeylonInstitution} for purchasable squares.
        \item \textbf{mortgageStatus}: \texttt{cannotMortgage} for non-purchasable, \texttt{noMortgage} for purchasable.
        \item \textbf{mortgageValue}: Half the current value (for properties, railways, utilities).
        \item \textbf{curruntValue}: Current market value of the square.
    \end{itemize}
    \item For \textbf{property squares}: additionally set \texttt{initialPrice}, \texttt{propertyGroup}, \texttt{baseRental}, \texttt{houseConstructionCost}, \texttt{hotelConstructionCost}, and \texttt{insuranceStatus = notInsured}.
    \item For \textbf{railway squares}: set rent tiers (250, 500, 1000, 2000) based on number owned.
    \item For \textbf{utility squares}: set \texttt{BaseRentalOfUtility = 999} (placeholder).
    \item For \textbf{special squares}: set \texttt{specililtyOfSquare} to \texttt{gotoJail}, \texttt{JailOrVisiting}, or \texttt{FreeParking}.
    \item For \textbf{tax squares}: set \texttt{taxAmount = 999}.
\end{enumerate}

\textbf{Board layout summary:}
\begin{table}[h]
\centering
\begin{tabular}{lll}
\toprule
Square & Name & Type \\
\midrule
0 & GO & go \\
1 & Pettah & property (brown) \\
2 & Community Development Fund & event \\
3 & Maradana & property (brown) \\
4 & Income Tax & tax \\
5 & Colombo Fort Railway Station & railway \\
6 & Bambalapitiya & property (light blue) \\
10 & Jail / Just Visiting & special \\
12 & Ceylon Electricity Board & utility \\
15 & Kandy Railway Station & railway \\
17 & Sri Lanka Insurance & insure \\
20 & Free Parking & special \\
25 & Galle Railway Station & railway \\
28 & National Water Supply & utility \\
30 & Go To Jail & special \\
35 & Jaffna Railway Station & railway \\
38 & Bank of Ceylon & bank \\
39 & Galle Face & property (dark blue) \\
\bottomrule
\end{tabular}
\caption{Selected board squares}
\end{table}

\newpage

% ============================================================
% SECTION 5: TURN ORDER AND DICE
% ============================================================
\section{Turn Order and Dice (\texttt{orderDetermine.c})}

\subsection{\texttt{dice\_roller(void)}}
Generates a random dice roll (two six-sided dice).

\textbf{Step-by-step logic:}
\begin{enumerate}[noitemsep]
    \item Roll die 1: \texttt{rand() \% 6 + 1} (produces 1--6).
    \item Roll die 2: \texttt{rand() \% 6 + 1}.
    \item If both dice are equal, set \texttt{doublesRolled = true}; otherwise \texttt{false}.
    \item Sum both dice into \texttt{rollValue}.
    \item Return the \texttt{diceRollType} struct.
\end{enumerate}

\subsection{\texttt{ranker(player*, player*, player*, player*)}}
Determines turn order by having all four players roll dice.

\textbf{Step-by-step logic:}
\begin{enumerate}[noitemsep]
    \item Store all four players in an array.
    \item Each player rolls via \texttt{dice\_roller()} and their \texttt{diceRoll} is recorded.
    \item Sort players in descending order of dice roll using bubble sort.
    \item Detect ties: if two or more consecutive players have the same roll, those players re-roll until all values are distinct.
    \item Re-sort the tied subgroup after re-rolling.
    \item Assign ranks 1--4 based on sorted order.
\end{enumerate}

\subsection{\texttt{finalRankAssign(player*, player*, player*, player*, player*)}}
Copies ranked players into their final ordered positions.

\textbf{Logic:}
\begin{enumerate}[noitemsep]
    \item If a player's rank is 1, copy them to \texttt{player\_1}.
    \item If rank is 2, copy to \texttt{player\_2}.
    \item If rank is 3, copy to \texttt{player\_3}.
    \item If rank is 4, copy to \texttt{player\_4}.
\end{enumerate}

\subsection{\texttt{roundCounter(context*, player*, player*, player*, player*)}}
Manages round advancement and prints round summaries.

\textbf{Step-by-step logic:}
\begin{enumerate}[noitemsep]
    \item Count active (non-bankrupt) players.
    \item Find the minimum number of completed laps among active players.
    \item If \texttt{completedLaps + 1 > currentBoardRound}, update the round counter.
    \item Print a detailed round summary for each non-bankrupt player: cash, net worth, properties, hotels, outstanding loan.
    \item Call \texttt{printMarketConditions()} to display economic state.
\end{enumerate}

\newpage

% ============================================================
% SECTION 6: PLAYER MOVEMENT
% ============================================================
\section{Player Movement (\texttt{moving.c})}

\subsection{\texttt{move(player*, square*, context*, playerPointers*)}}
Handles a player's turn: rolling dice, moving, and resolving the landed square.

\textbf{Step-by-step logic:}
\begin{enumerate}[noitemsep]
    \item Call \texttt{dice\_roller()} to get the roll result.
    \item If the player is \textbf{in jail}:
    \begin{itemize}[noitemsep]
        \item Increment \texttt{jailRoundCounter}.
        \item Call \texttt{jailLogic()} to check if the player can leave jail.
        \item Return immediately (no movement this turn).
    \end{itemize}
    \item Print the roll value.
    \item Calculate new position: \texttt{(currentSquare + rollValue) \% 40} (wraps around the board).
    \item Update \texttt{totalsteps}.
    \item Print movement from old position to new position.
    \item Call \texttt{resolveSquare()} to handle the landed square's effect.
    \item \textbf{Pass-Go bonus:} If the player's new position is less than the old position (and not on square 10 or 0), and the player has no debt, award LKR 2,000.
    \item Update \texttt{laps = totalsteps / 40}.
\end{enumerate}

\newpage

% ============================================================
% SECTION 7: SQUARE RESOLUTION
% ============================================================
\section{Square Resolution (\texttt{resolveSquare.c})}

\subsection{\texttt{resolveSquare(player*, square*, context*, playerPointers*)}}
The main dispatcher that routes square resolution to the appropriate handler.

\textbf{Step-by-step logic:}
\begin{enumerate}[noitemsep]
    \item \textbf{Debt timeout check:} If the player has outstanding debt and 20 rounds have passed since taking it, confiscate all assets and auction them off. Mark the player as debt-free.
    \item Read the \texttt{squareType} of the current square.
    \item Switch on the type and call the corresponding resolver:
    \begin{itemize}[noitemsep]
        \item \texttt{go} $\rightarrow$ \texttt{resolveGO()}
        \item \texttt{special} $\rightarrow$ \texttt{resolveSpecial()}
        \item \texttt{railway} $\rightarrow$ \texttt{resolveRailway()}
        \item \texttt{utility} $\rightarrow$ \texttt{resolveUtility()}
        \item \texttt{event} $\rightarrow$ \texttt{resolveEvent()}
        \item \texttt{insure} $\rightarrow$ \texttt{resolveInsure()}
        \item \texttt{tax} $\rightarrow$ \texttt{resolveTax()}
        \item \texttt{bank} $\rightarrow$ \texttt{resolveBank()}
        \item \texttt{property} $\rightarrow$ \texttt{resolveProperty()}
    \end{itemize}
\end{enumerate}

\subsection{\texttt{jailLogic(player*, bool doublesRolled)}}
Determines whether a player can leave jail.

\textbf{Logic (conservative banker):}
\begin{enumerate}[noitemsep]
    \item If doubles rolled $\rightarrow$ release immediately.
    \item Else if jailed for $>3$ rounds $\rightarrow$ release automatically.
\end{enumerate}

\textbf{Logic (all other players):}
\begin{enumerate}[noitemsep]
    \item If cash $> 300$ $\rightarrow$ pay LKR 300 bailout and release.
    \item Else if jailed for $>3$ rounds $\rightarrow$ release automatically.
    \item Else if doubles rolled $\rightarrow$ release immediately.
\end{enumerate}

\subsection{\texttt{resolveGO(player*)}}
Adds LKR 2,000 to the player's cash for landing on GO.

\subsection{\texttt{resolveSpecial(player*, square*)}}
Handles special squares (Go To Jail, Free Parking).
\begin{enumerate}[noitemsep]
    \item If the square is \texttt{gotoJail}: move player to square 10, set jail status to \texttt{inside}, reset jail counter.
    \item Otherwise (Free Parking, Jail/Visiting): no action.
\end{enumerate}

\subsection{\texttt{resolveRailway(player*, square*, playerPointers*, context*)}}
Handles railway squares (buying and rent payment).

\textbf{If owned by bank:}
\begin{enumerate}[noitemsep]
    \item Check the player type and call the corresponding buy condition:
    \begin{itemize}[noitemsep]
        \item \textbf{Aggressive Investor}: buys if cash after purchase $\geq$ 2000.
        \item \textbf{Risk Taker}: buys if cash $\geq$ current value.
        \item \textbf{Conservative Banker}: buys if remaining cash $> 30\%$ of total cash.
        \item \textbf{Opportunistic Trader}: buys if cash $\geq$ value AND not in Economic Recession.
    \end{itemize}
    \item If condition is met, transfer ownership and increment \texttt{noOfRailways}.
\end{enumerate}

\textbf{If owned by another player:}
\begin{enumerate}[noitemsep]
    \item Look up the rent tier based on how many railways the owner has (1--4).
    \item Deduct rent from the current player, add to the owner.
\end{enumerate}

\subsection{\texttt{resolveProperty(player*, square*, context*, playerPointers*)}}
The most complex resolver; handles buying, rent, and building.

\textbf{Case 1: Property owned by bank (buying decision):}
\begin{enumerate}[noitemsep]
    \item \textbf{Aggressive Investor}: Buys if cash $> 1000 +$ initial price.
    \item \textbf{Conservative Banker}: Buys if remaining cash $> 50\%$ of total AND not in recession.
    \item \textbf{Risk Taker}: Buys if cash $>$ initial price.
    \item \textbf{Opportunistic Trader}: Buys if cash $>$ initial price AND not in recession AND inflation $> 0$.
\end{enumerate}

\textbf{Case 2: Property owned by another player (rent):}
\begin{enumerate}[noitemsep]
    \item If cash $>$ rent AND property not mortgaged $\rightarrow$ pay rent via \texttt{payRent()}.
    \item If cash $<$ rent AND no properties owned $\rightarrow$ mark bankrupt.
    \item If cash $<$ rent BUT has properties $\rightarrow$ trigger \texttt{AggrNoCashAuction()} to sell properties until rent can be paid.
\end{enumerate}

\textbf{Case 3: Property owned by the player (building):}
\begin{enumerate}[noitemsep]
    \item Check monopoly eligibility via \texttt{checkForMonopoly()}.
    \item \textbf{Aggressive Investor}: builds if has monopoly and can afford house cost.
    \item \textbf{Conservative Banker}: builds if has monopoly, remaining cash $> 50\%$ of total, and not in recession.
    \item \textbf{Risk Taker}: builds if has monopoly and can afford house cost.
    \item \textbf{Opportunistic Trader}: builds if has monopoly, can afford cost, and inflation $\leq 0$.
    \item \textbf{Building logic:} If $<4$ houses, build a house. If 4 houses and can afford hotel, upgrade to hotel.
    \item \textbf{Rent scaling:}
    \begin{itemize}[noitemsep]
        \item 1 house: $2\times$ base rent
        \item 2 houses: $3\times$ base rent
        \item 3 houses: $5\times$ base rent
        \item 4 houses: $7\times$ base rent
        \item Hotel: $10\times$ base rent
    \end{itemize}
\end{enumerate}

\subsection{\texttt{checkForMonopoly(player*, square*)}}
Determines if a player can build on a property.

\textbf{Step-by-step logic:}
\begin{enumerate}[noitemsep]
    \item Get the property group of the current square.
    \item Check that the player owns ALL properties in that group (set \texttt{ownsAMonopoly}).
    \item Enforce even building: all properties in the group must have the same effective house count (houses + hotel counted as 5).
    \item Check that no property in the group is mortgaged.
    \item Return \texttt{true} only if all three conditions are met.
\end{enumerate}

\newpage

% ============================================================
% SECTION 8: FINANCE
% ============================================================
\section{Finance and Economics (\texttt{finance.c})}

\subsection{\texttt{networthEvaluate(player*, player*, player*, player*, square*)}}
Recalculates net worth for all four players each round.

\textbf{Step-by-step logic:}
\begin{enumerate}[noitemsep]
    \item For each player:
    \begin{enumerate}[noitemsep]
        \item Iterate over all 40 squares.
        \item For each property owned by the player:
        \begin{itemize}[noitemsep]
            \item If mortgaged: value = mortgage value.
            \item Otherwise: value = current value + (houses $\times$ house cost) + (hotels $\times$ hotel cost).
        \end{itemize}
        \item For railways/utilities: value = current value (or mortgage value if mortgaged).
        \item Sum all asset values.
    \end{enumerate}
    \item Calculate: \texttt{netWorth = cash + totalAssetValue - outstandingLoan}.
    \item Calculate: \texttt{MaxEligibleLoanAmount = totalMortgageValues $\times$ 0.75}.
\end{enumerate}

\subsection{\texttt{resolveBank(player*, square*, context*)}}
Handles loan borrowing and repayment at the Bank of Ceylon.

\textbf{If player has no debt (borrowing):}
\begin{enumerate}[noitemsep]
    \item \textbf{Aggressive Investor}: Borrows if has a monopoly AND cash $\leq$ 5000.
    \item \textbf{Conservative Banker}: Borrows if net worth $< 5000$.
    \item \textbf{Risk Taker}: Always borrows the maximum eligible amount.
    \item \textbf{Opportunistic Trader}: Borrows only during Government Housing Programme or Stock Market Boom events.
\end{enumerate}

\textbf{If player has debt (repayment):}
\begin{enumerate}[noitemsep]
    \item All player types repay if cash $> 2 \times$ outstanding loan.
    \item Repayment clears the debt entirely.
\end{enumerate}

\subsection{\texttt{playerHasaMonopoly(player*, square*)}}
Checks if a player owns all properties in any single color group.

\textbf{Logic:}
\begin{enumerate}[noitemsep]
    \item Initialize 8 boolean flags (one per color group), all \texttt{true}.
    \item Scan all 40 squares. For each property, if the current player does not own it, set that group's flag to \texttt{false}.
    \item Return \texttt{true} if ANY group flag is still \texttt{true}.
\end{enumerate}

\subsection{\texttt{payRent(player*, square*)}}
Handles rent transfer between players.

\textbf{Logic:}
\begin{enumerate}[noitemsep]
    \item If player's cash $<$ rent amount $\rightarrow$ return \texttt{false}.
    \item Deduct rent from the player's cash.
    \item Add rent to the property owner's cash.
    \item Print the transaction.
    \item Return \texttt{true}.
\end{enumerate}

\subsection{\texttt{bankruptCheck(playerPointers*, int*)}}
Checks all players for bankruptcy.

\textbf{Logic:}
\begin{enumerate}[noitemsep]
    \item Reset \texttt{noOfBankruptPlayers} to 0.
    \item For each player: if net worth $\leq 0$, mark as bankrupt (if not already).
    \item Count total bankrupt players.
\end{enumerate}

\subsection{\texttt{sellingAuction(...)}} and \texttt{bankruptAuction(...)}
Simulate property auctions with player-specific bidding strategies.

\textbf{Auction logic (sellingAuction):}
\begin{enumerate}[noitemsep]
    \item Set starting price to half the property's current value.
    \item Collect the 3 non-owner players as bidders.
    \item Loop until a winner is determined:
    \begin{enumerate}[noitemsep]
        \item If the owner has no properties left $\rightarrow$ bankrupt.
        \item Each bidder decides whether to bid based on their player type:
        \begin{itemize}[noitemsep]
            \item \textbf{Aggressive Investor}: Bids if cash $\geq$ highest bid + 250 AND bid $< 120\%$ of current value.
            \item \textbf{Risk Taker}: Bids if cash $\geq$ highest bid + 250 (no upper limit).
            \item \textbf{Conservative Banker}: Bids if remaining cash $> 50\%$ AND bid $< 133\%$ of current value.
            \item \textbf{Opportunistic Trader}: Bids if cash $>$ highest bid + 250 AND bid $< 4000$.
        \end{itemize}
        \item If a bidder outbids the current highest bidder, they become the new highest bidder.
        \item If no one outbids, the property sells at mortgage value to the bank.
    \end{enumerate}
    \item Winner pays, receives the property (houses/hotels reset to 0).
\end{enumerate}

\textbf{Bankrupt auction (bankruptAuction):} Same logic but the property reverts to the bank if no one bids.

\subsection{\texttt{AggrNoCashAuction(square*, player*, playerPointers*, context*)}}
Handles the case where a player cannot pay rent.

\textbf{Logic:}
\begin{enumerate}[noitemsep]
    \item Loop while cash $<$ rent AND player has properties:
    \begin{enumerate}[noitemsep]
        \item Recalculate net worth.
        \item Iterate through each color group (brown through dark blue).
        \item For each property owned by the player in that group, trigger \texttt{sellingAuction()}.
        \item After each sale, attempt to pay rent via \texttt{payRent()}.
        \item If rent is paid, return immediately.
    \end{enumerate}
    \item If the loop exits and rent still cannot be paid $\rightarrow$ mark bankrupt.
\end{enumerate}

\newpage

% ============================================================
% SECTION 9: ECONOMIC EVENTS
% ============================================================
\section{Economic Events (\texttt{events.c})}

\subsection{\texttt{econEventActivate(square*, context*)}}
Triggers a random economic event every 15 rounds.

\textbf{Step-by-step logic:}
\begin{enumerate}[noitemsep]
    \item Define the array of 8 economic event types.
    \item If an event is currently active, \textbf{deactivate} it (call the corresponding \texttt{\_Deactivate} function).
    \item Select a random event: \texttt{rand() \% 8}.
    \item Store it in \texttt{currentActiveEconEvent}.
    \item Call the corresponding activation function.
    \item Record the round this event happened.
\end{enumerate}

\subsection{\texttt{inflationRateRelease(square*, context*)}}
Applies inflation every 10 rounds.

\textbf{Step-by-step logic:}
\begin{enumerate}[noitemsep]
    \item Select a random inflation rate from $\{-3, 0, 2, 5, 8, 12\}$.
    \item Store in \texttt{currentInflation}.
    \item For all 40 squares, adjust:
    \begin{itemize}[noitemsep]
        \item \texttt{curruntValue} $\times (1 + \textrm{inflation}/100)$
        \item \texttt{currentRentalofProperty} $\times (1 + \textrm{inflation}/100)$
        \item \texttt{initialPrice} $\times (1 + \textrm{inflation}/100)$
    \end{itemize}
    \item Record the round this inflation happened.
\end{enumerate}

\subsection{\texttt{printMarketConditions(context*)}}
Displays the current economic state.

\textbf{Logic:}
\begin{enumerate}[noitemsep]
    \item Calculate rounds remaining for inflation and economic event.
    \item Print inflation rate (with sign).
    \item Print the active economic event description via a switch statement.
    \item Print rounds remaining for the event.
\end{enumerate}

\newpage

% ============================================================
% SECTION 10: ECONOMIC EVENT EFFECTS
% ============================================================
\section{Economic Event Effects (\texttt{national\_event\_functions.c})}

Each economic event has an \texttt{Activate} and \texttt{Deactivate} function that modifies board properties.

\begin{table}[h]
\centering
\small
\begin{tabular}{p{3.5cm}p{5cm}p{5cm}}
\toprule
\textbf{Event} & \textbf{Activation Effect} & \textbf{Deactivation Effect} \\
\midrule
Tourism Boom & Hotel rents $\times 2$; Yellow group value $\times 1.15$ & Reverse both \\
Fuel Crisis & Railway rents $\times 2$; House \& hotel costs $\times 1.2$ & Reverse both \\
Heavy Monsoon & Yellow group value $\times 0.9$ & Divide by 0.9 \\
Economic Recession & Property value $\times 0.85$; Rents $\times 0.9$ & Reverse both \\
Stock Market Boom & Property value $\times 1.1$ & Divide by 1.1 \\
Gov.\ Housing Programme & House construction cost $\times 0.75$ & Divide by 0.75 \\
Foreign Investment & Orange \& Red group value $\times 1.2$ & Divide by 1.2 \\
Political Unrest & Hotel rents $\times 0.5$ & Divide by 0.5 \\
\bottomrule
\end{tabular}
\caption{Economic event effects}
\end{table}

\textbf{Pattern:} Each activation modifies the board, and the corresponding deactivation restores the original values by applying the inverse operation.

\newpage

% ============================================================
% SECTION 11: PLAYER BEHAVIORS
% ============================================================
\section{Player Behaviors (\texttt{players.c})}

\subsection{Railway Buying Conditions}

Each player type has a different decision criterion for buying railways:

\begin{table}[h]
\centering
\begin{tabular}{ll}
\toprule
\textbf{Player Type} & \textbf{Buy Condition} \\
\midrule
Aggressive Investor & Cash $-$ price $\geq 2000$ \\
Risk Taker & Cash $\geq$ price \\
Conservative Banker & Cash $-$ price $> 30\%$ of cash \\
Opportunistic Trader & Cash $\geq$ price AND not in recession \\
\bottomrule
\end{tabular}
\caption{Railway buying conditions}
\end{table}

\subsection{Property Buying Conditions (from resolveProperty)}

\begin{table}[h]
\centering
\small
\begin{tabular}{p{3.5cm}p{8cm}}
\toprule
\textbf{Player Type} & \textbf{Buy Condition} \\
\midrule
Aggressive Investor & Cash $> 1000 +$ initial price \\
Conservative Banker & Remaining cash $> 50\%$ of total AND not in recession \\
Risk Taker & Cash $>$ initial price \\
Opportunistic Trader & Cash $>$ price AND not in recession AND inflation $> 0$ \\
\bottomrule
\end{tabular}
\caption{Property buying conditions}
\end{table}

\subsection{Building Conditions (from resolveProperty)}

\begin{table}[h]
\centering
\small
\begin{tabular}{p{3.5cm}p{8cm}}
\toprule
\textbf{Player Type} & \textbf{Build Condition} \\
\midrule
Aggressive Investor & Has monopoly AND can afford house cost \\
Conservative Banker & Has monopoly AND remaining cash $> 50\%$ AND not in recession \\
Risk Taker & Has monopoly AND can afford house cost \\
Opportunistic Trader & Has monopoly AND can afford cost AND inflation $\leq 0$ \\
\bottomrule
\end{tabular}
\caption{Building conditions}
\end{table}

\newpage

% ============================================================
% SECTION 12: OVERALL FLOW
% ============================================================
\section{Overall Program Flow}

\begin{enumerate}
    \item \texttt{main()} calls \texttt{startGame()}.
    \item \texttt{startGame()} initializes players, board, and turn order.
    \item Main loop runs up to 500 rounds:
    \begin{enumerate}[noitemsep]
        \item Each non-bankrupt player takes a turn via \texttt{move()}.
        \item \texttt{move()} rolls dice, updates position, calls \texttt{resolveSquare()}.
        \item \texttt{resolveSquare()} dispatches to type-specific handlers.
        \item Properties may be bought, rent paid, or buildings constructed.
        \item Net worths are recalculated each round.
        \item Bankrupt players are identified.
        \item Economic events trigger every 15 rounds.
        \item Inflation adjusts every 10 rounds.
    \end{enumerate}
    \item Game ends when 3 players go bankrupt or 500 rounds pass.
    \item Winner is declared based on highest net worth or being the sole survivor.
\end{enumerate}

\begin{figure}[h]
\centering
\begin{tcolorbox}[colback=white, colframe=black, width=0.85\textwidth]
\centering
\texttt{main()} $\rightarrow$ \texttt{startGame()} \\
$\downarrow$ \\
Initialize players, board, turn order \\
$\downarrow$ \\
\textbf{Main Loop} (up to 500 rounds) \\
$\downarrow$ \\
\texttt{move()} $\rightarrow$ \texttt{dice\_roller()} $\rightarrow$ \texttt{resolveSquare()} \\
$\downarrow$ \\
\texttt{resolveProperty()} / \texttt{resolveRailway()} / etc. \\
$\downarrow$ \\
\texttt{networthEvaluate()} $\rightarrow$ \texttt{bankruptCheck()} \\
$\downarrow$ \\
\texttt{econEventActivate()} / \texttt{inflationRateRelease()} \\
$\downarrow$ \\
Declare Winner
\end{tcolorbox}
\caption{High-level program flow}
\end{figure}

\end{document}
